% Preámbulo compartido para todas las clases de Beamer.
% Cada clase lo incluye con % Preámbulo compartido para todas las clases de Beamer.
% Cada clase lo incluye con % Preámbulo compartido para todas las clases de Beamer.
% Cada clase lo incluye con \input{../../latex/preamble.tex}

\documentclass[aspectratio=169,11pt]{beamer}

\usetheme{Madrid}
\setbeamertemplate{navigation symbols}{}
\setbeamertemplate{footline}[frame number]

\usepackage[utf8]{inputenc}
\usepackage[spanish]{babel}
\usepackage{amsmath,amssymb}
\usepackage{graphicx}
\usepackage{booktabs}
\usepackage{listings}
\usepackage{xcolor}

\lstset{
  basicstyle=\ttfamily\small,
  keywordstyle=\color{blue},
  commentstyle=\color{gray},
  stringstyle=\color{teal},
  showstringspaces=false,
  breaklines=true,
  frame=single
}

\newcommand{\definicion}[1]{%
  \begin{block}{Definición}
    #1
  \end{block}
}

\newcommand{\ejemplo}[1]{%
  \begin{exampleblock}{Ejemplo}
    #1
  \end{exampleblock}
}

\newcommand{\importante}[1]{%
  \begin{alertblock}{Importante}
    #1
  \end{alertblock}
}

\author{Agustín Gurvich}
\institute{Introducción a Machine Learning}


\documentclass[aspectratio=169,11pt]{beamer}

\usetheme{Madrid}
\setbeamertemplate{navigation symbols}{}
\setbeamertemplate{footline}[frame number]

\usepackage[utf8]{inputenc}
\usepackage[spanish]{babel}
\usepackage{amsmath,amssymb}
\usepackage{graphicx}
\usepackage{booktabs}
\usepackage{listings}
\usepackage{xcolor}

\lstset{
  basicstyle=\ttfamily\small,
  keywordstyle=\color{blue},
  commentstyle=\color{gray},
  stringstyle=\color{teal},
  showstringspaces=false,
  breaklines=true,
  frame=single
}

\newcommand{\definicion}[1]{%
  \begin{block}{Definición}
    #1
  \end{block}
}

\newcommand{\ejemplo}[1]{%
  \begin{exampleblock}{Ejemplo}
    #1
  \end{exampleblock}
}

\newcommand{\importante}[1]{%
  \begin{alertblock}{Importante}
    #1
  \end{alertblock}
}

\author{Agustín Gurvich}
\institute{Introducción a Machine Learning}


\documentclass[aspectratio=169,11pt]{beamer}

\usetheme{Madrid}
\setbeamertemplate{navigation symbols}{}
\setbeamertemplate{footline}[frame number]

\usepackage[utf8]{inputenc}
\usepackage[spanish]{babel}
\usepackage{amsmath,amssymb}
\usepackage{graphicx}
\usepackage{booktabs}
\usepackage{listings}
\usepackage{xcolor}

\lstset{
  basicstyle=\ttfamily\small,
  keywordstyle=\color{blue},
  commentstyle=\color{gray},
  stringstyle=\color{teal},
  showstringspaces=false,
  breaklines=true,
  frame=single
}

\newcommand{\definicion}[1]{%
  \begin{block}{Definición}
    #1
  \end{block}
}

\newcommand{\ejemplo}[1]{%
  \begin{exampleblock}{Ejemplo}
    #1
  \end{exampleblock}
}

\newcommand{\importante}[1]{%
  \begin{alertblock}{Importante}
    #1
  \end{alertblock}
}

\author{Agustín Gurvich}
\institute{Introducción a Machine Learning}


\title{Python Científico y Preprocesamiento de Datos}
\subtitle{numpy, pandas, matplotlib y preparar datos para un modelo}
\date{}

\begin{document}

\frame{\titlepage}


\part{Python científico}
\frame{\partpage}

\section{¿Por qué no alcanza con Python puro?}

\begin{frame}{¿Por qué no alcanza con Python puro?}
  \begin{itemize}
    \item Las listas de Python no fueron pensadas para hacer cuentas con miles o millones de números
    \item Para trabajar con datos en serio usamos tres librerías:
  \end{itemize}
  \begin{itemize}
    \item \textbf{numpy}: vectores y matrices de números
    \item \textbf{pandas}: tablas (filas y columnas)
    \item \textbf{matplotlib}: gráficos
  \end{itemize}
  \pause
  \importante{
    Casi toda librería de ML en Python, incluyendo scikit-learn, está
    construida arriba de numpy y pandas.
  }
\end{frame}

\section{numpy}

\begin{frame}[fragile]{numpy: arrays}
  \definicion{
    Un array de numpy es una lista de números que se puede operar
    directamente, sin escribir loops.
  }
  \begin{lstlisting}[language=Python]
import numpy as np

x = np.array([1, 2, 3, 4])
x * 2
# array([2, 4, 6, 8])
  \end{lstlisting}
\end{frame}

\begin{frame}[fragile]{numpy: operaciones sin loops}
  \begin{lstlisting}[language=Python]
edades = np.array([15, 16, 17, 18, 17])

edades.mean()   # 16.6
edades.std()    # desvio estandar
edades.max()    # 18
edades[edades > 16]   # array([17, 18, 17])
  \end{lstlisting}
  \pause
  \importante{
    Nada de \texttt{for}. Estas operaciones se aplican a todo el array de una,
    y son mucho más rápidas que un loop escrito a mano.
  }
\end{frame}

\section{pandas}

\begin{frame}[fragile]{pandas: DataFrames}
  \definicion{
    Un DataFrame es una tabla: filas y columnas, con nombres, donde cada
    columna puede tener un tipo de dato distinto (números, texto, fechas).
  }
  \begin{lstlisting}[language=Python]
import pandas as pd

df = pd.read_csv("datos.csv")
df.head()   # primeras 5 filas
  \end{lstlisting}
\end{frame}

\begin{frame}[fragile]{pandas: elegir y filtrar}
  \begin{lstlisting}[language=Python]
df["edad"]                    # una columna
df[["edad", "nota"]]          # varias columnas

df[df["edad"] > 17]           # filas que cumplen una condicion
  \end{lstlisting}
\end{frame}

\begin{frame}[fragile]{pandas: funciones que van a usar seguido}
  \begin{itemize}
    \item \texttt{df.info()}: tipos de datos y valores faltantes
    \item \texttt{df.describe()}: media, desvío, mínimo, máximo de cada columna
    \item \texttt{df.groupby("columna")}: agrupar filas por un valor
    \item \texttt{df.sort\_values("columna")}: ordenar
  \end{itemize}
\end{frame}

\section{matplotlib}

\begin{frame}[fragile]{matplotlib: visualizar}
  \begin{lstlisting}[language=Python]
import matplotlib.pyplot as plt

fig, ax = plt.subplots()
ax.scatter(df["edad"], df["nota"])
plt.show()
  \end{lstlisting}
  \pause
  \texttt{fig} es el lienzo entero, \texttt{ax} es donde se dibuja. Ya lo
  probamos con el dataset Iris.
\end{frame}

\begin{frame}{Resumen: numpy, pandas, matplotlib}
  \begin{itemize}
    \item numpy: cuentas rápidas sobre números, sin loops
    \item pandas: tablas, filtrar, agrupar, resumir
    \item matplotlib: ver los datos antes y después de modelar
  \end{itemize}
\end{frame}

\part{Preprocesamiento de datos}
\frame{\partpage}

\section{De datos crudos a datos usables}

\begin{frame}{De datos a datos usables}
  \importante{
    Ningún modelo entrena bien con datos sucios. En un proyecto real,
    esta parte suele llevar más tiempo que elegir y entrenar el modelo.
  }
  Antes de entrenar cualquier cosa hay que resolver al menos tres
  problemas típicos:
  \begin{enumerate}
    \item Datos faltantes
    \item Variables que no son números (categóricas)
    \item Variables en escalas muy distintas
  \end{enumerate}
\end{frame}

\section{Datos faltantes}

\begin{frame}[fragile]{Datos faltantes}
  \definicion{
    Un dato faltante (\texttt{NaN}) es un valor que no se registró:
    una encuesta sin responder, un sensor que falló, etc.
  }
  \begin{lstlisting}[language=Python]
df.isna().sum()          # cuantos faltan por columna

df.dropna()              # opcion 1: eliminar filas con faltantes
df.fillna(df.mean())     # opcion 2: completar con la media
  \end{lstlisting}
\end{frame}

\begin{frame}{Datos faltantes: ¿eliminar o completar?}
  \begin{itemize}
    \item Eliminar filas: fácil, pero se pierde información si hay muchos faltantes
    \item Completar con la media/mediana: mantiene el tamaño del dataset, pero "inventa" un valor
  \end{itemize}
  \pause
  No hay una respuesta única, depende de cuántos datos falten y de qué
  tan importante sea esa columna.
\end{frame}

\section{Variables categóricas}

\begin{frame}{Variables categóricas}
  \definicion{
    Columnas con categorías en vez de números: color, ciudad, tipo de
    producto.
  }
  \importante{
    Los modelos de ML trabajan con números. Una columna de texto hay que
    convertirla antes de entrenar.
  }
\end{frame}

\begin{frame}[fragile]{One hot encoding}
  Convierte cada categoría en una columna nueva de ceros y unos.
  \begin{lstlisting}[language=Python]
pd.get_dummies(df, columns=["color"])
  \end{lstlisting}
  \begin{center}
  \begin{tabular}{lccc}
    \toprule
    color & color\_rojo & color\_verde & color\_azul \\
    \midrule
    rojo & 1 & 0 & 0 \\
    verde & 0 & 1 & 0 \\
    azul & 0 & 0 & 1 \\
    \bottomrule
  \end{tabular}
  \end{center}
\end{frame}

\section{Escalado de variables}

\begin{frame}{Escalado de variables}
  \ejemplo{
    Un dataset con \texttt{edad} (0 a 100) y \texttt{sueldo} (0 a
    1.000.000). Para el modelo, una diferencia de sueldo de 10.000 "pesa"
    mucho más que una diferencia de edad de 10 años, solo porque el
    número es más grande.
  }
  \pause
  \definicion{
    Escalar significa transformar las variables numéricas para que
    queden en rangos comparables entre sí.
  }
\end{frame}

\begin{frame}[fragile]{Escalado: StandardScaler}
  \begin{lstlisting}[language=Python]
from sklearn.preprocessing import StandardScaler

scaler = StandardScaler()
X_escalado = scaler.fit_transform(X)
  \end{lstlisting}
  Cada columna queda con media 0 y desvío estándar 1. Ninguna variable
  domina a las demás solo por su escala original.
\end{frame}

\section{Train / test split}

\begin{frame}{Train / test split}
  \definicion{
    Separar el dataset en dos partes: una para entrenar el modelo
    (\textit{train}) y otra que el modelo nunca ve durante el
    entrenamiento (\textit{test}), para evaluarlo después.
  }
  \pause
  \importante{
    Evaluar un modelo con los mismos datos que usó para entrenar es como
    corregir un examen dándole las respuestas al alumno antes de
    tomarlo.
  }
\end{frame}

\begin{frame}[fragile]{Train / test split en código}
  \begin{lstlisting}[language=Python]
from sklearn.model_selection import train_test_split

X_train, X_test, y_train, y_test = train_test_split(
    X, y, test_size=0.2, random_state=0
)
  \end{lstlisting}
  Una división típica: 70\% para entrenar, 30\% para evaluar.
\end{frame}

\begin{frame}{Resumen}
  \begin{itemize}
    \item numpy, pandas y matplotlib son la base para manipular y ver datos
    \item Faltantes, categóricas y escalado hay que resolverlos antes de entrenar
    \item El test set se separa antes de entrenar y no se toca hasta evaluar
  \end{itemize}
  \vspace{1em}
  \centering
  \Large ¿Preguntas?
\end{frame}

\end{document}
